%% ======================================================================
%% preamble-figures.tex
%% ----------------------------------------------------------------------
%% Every definition the figures need, kept identical to main.tex's preamble.
%% Shared by the standalone builds in this folder; main.tex carries its own
%% copy so that the submitted source stays a single .tex file.
%% ======================================================================

\providecommand{\method}{\textsc{ESV}}

%% --- comparison-table glyphs (fig_trap uses \yes / \no) ---------------
\definecolor{okgreen}{RGB}{0,122,87}
\definecolor{okred}{RGB}{176,42,32}
\providecommand{\yes}{\textcolor{okgreen}{\ding{51}}}
\providecommand{\no}{\textcolor{okred}{\ding{55}}}
\providecommand{\na}{\textcolor{black!45}{--}}

%% --- figure fills, shaded by what each step costs ---------------------
\definecolor{cshared}{RGB}{226,226,226}
\definecolor{ccode}{RGB}{205,228,243}
\definecolor{ccall}{RGB}{252,230,192}
\definecolor{ccallb}{RGB}{191,118,0}
\definecolor{cout}{RGB}{206,238,227}
\definecolor{ctrap}{RGB}{250,222,218}

%% --- plot colours and markers -----------------------------------------
%% Marker shapes are core TikZ nodes, so every glyph is vector and uses the
%% document's Type 1 text fonts: no Type 3 subsets, no raster, nothing to crop.
\definecolor{pblue}{RGB}{31,97,160}
\definecolor{pgreen}{RGB}{0,122,87}
\definecolor{pgrey}{RGB}{95,95,95}
\providecommand{\mkopensq}[2]{\node[draw=#2,fill=white,line width=.6pt,rectangle,inner sep=1.5pt] at #1 {};}
\providecommand{\mkfillsq}[2]{\node[draw=#2,fill=#2,rectangle,inner sep=1.5pt] at #1 {};}
\providecommand{\mkopenci}[2]{\node[draw=#2,fill=white,line width=.6pt,circle,inner sep=1.5pt] at #1 {};}
\providecommand{\mkfillci}[2]{\node[draw=#2,fill=#2,circle,inner sep=1.5pt] at #1 {};}
%% #1 x, #2 centre, #3 half-width of the interval, #4 colour (needs x=1cm)
\providecommand{\ebar}[4]{%
  \draw[#4,line width=.5pt] (#1,{#2-#3}) -- (#1,{#2+#3});
  \draw[#4,line width=.5pt] ({#1-0.055},{#2-#3}) -- ({#1+0.055},{#2-#3});
  \draw[#4,line width=.5pt] ({#1-0.055},{#2+#3}) -- ({#1+0.055},{#2+#3});}

%% --- \cost is used inside the three schematic figures -----------------
\providecommand{\cost}{\scriptsize}
